\par\noindent\textbf{(1)} 三维自由电子态密度g(ε)=C√ε（C含体积和两倍自旋），总粒子数N=∫₀∞g(ε)f\_F(ε,μ)dε。Sommerfeld展开给N=∫₀\textasciicircum{}μg(ε)dε+(π²/6)(k\_BT)²g'(μ)+O(T⁴)，而T=0时N=(2C/3)EF\textasciicircum{}\{3/2\}。
\par\noindent\textbf{(2)} 设μ=EF+δμ且δμ∝T²，粒子数方程一阶变分为0=g(EF)δμ+(π²/6)(k\_BT)²g'(EF)。因g'/g=1/(2EF)，得μ≈EF[1−(π²/12)(k\_BT/EF)²]，低温化学势略下降而非固定于EF。
\par\noindent\textbf{(3)} 能量U=∫εg(ε)f\_F dε，对H(ε)=εg(ε)=Cε\textasciicircum{}\{3/2\}用同样展开。代入μ的T²变化，得U≈(3/5)NEF+(π²/4)N(k\_BT)²/EF，所以定体积电子热容 C\_V=(∂U/∂T)\_N≈(π²/2)Nk\_B²T/EF。
\par\noindent\textbf{(4)} 只有费米面附近宽度约k\_BT的电子可以重新分布，活跃粒子比例约T/TF，故热容随T线性而非经典3Nk\_B/2。T接近TF、强相互作用或非抛物色散时，Sommerfeld的小参数展开失效，需计算具体态密度。
